% Paper_8_Bond_Sphere_Sturmian.tex — GeoVac Paper 8-9 (Merged)
% Bond Sphere Geometry + Sturmian Basis
% Status: Draft, March 2026
% Target: arXiv (quant-ph / physics.comp-ph)
% Merges former Paper 8 (Bond Sphere) and Paper 9 (Sturmian Basis)
\documentclass[aps,pra,twocolumn,superscriptaddress,longbibliography]{revtex4-2}

\usepackage{amsmath,amssymb,amsthm,graphicx,xcolor,bm,braket,dcolumn}

\newtheorem{theorem}{Theorem}
\newtheorem{corollary}{Corollary}

\begin{document}

\title{Bond Sphere: Diatomic Molecules as $S^3$ and\\
their Natural Sturmian Basis}
\thanks{Paper 8--9 in the Geometric Vacuum series}

\author{J.~Loutey}
\affiliation{Independent Researcher, Kent, Washington}

\date{March 8, 2026}

\begin{abstract}
A diatomic molecule is geometrically a single three-sphere $S^3$ with two
weighted poles---the \emph{bond sphere}---in which the internuclear distance
$R$ encodes as a polar angle $\gamma$ via
$\cos\gamma = (p_0^2 - p_R^2)/(p_0^2 + p_R^2)$ and all cross-atom matrix
elements are $SO(4)$ Wigner $D$-matrix elements evaluated at the bond
rotation $e^{i\gamma A_y}$.
\emph{Positive results:}
The bond sphere geometry is consistent, with the $\gamma(R)$ encoding,
heteronuclear egg metric, Runge--Lenz alignment, and kinetic membrane all
following from $SO(4)$ representation theory.
The $D$-matrix enforces selection rules for $\sigma$-bonds:
$D^2_{10,10}(\gamma) = 1$ and $D^2_{00,10}(\gamma) = 0$ for all $\gamma$.
\emph{Negative theorem:}
In the molecular Sturmian basis with shared momentum scale $p_0$,
$H_{ij} \propto S_{ij}$; eigenvalues are independent of $R$.
Binding requires the $R$-dependent prolate spheroidal spectrum $\beta_k(R)$,
reintroducing continuous geometry.
The resolution---discretizing the Laplace--Beltrami operator in prolate
spheroidal coordinates---is developed in the companion
Paper~11~\cite{loutey_paper11}.
\end{abstract}

\maketitle


%% ========================================================================
\section{Introduction}
\label{sec:intro}
%% ========================================================================

In the GeoVac framework~\cite{loutey_paper0,loutey_paper1,loutey_paper7},
quantum numbers $(n, l, m)$ are not abstract labels but geometric coordinates
on a three-sphere $S^3$.  The principal quantum number $n$ indexes latitude
shells, each containing $n^2$ nodes; the orbital angular momentum $l$ and
magnetic quantum number $m$ locate a node within its shell.  The graph
Laplacian over these nodes is mathematically equivalent to the Laplace--Beltrami
operator on the unit $S^3$, and the continuous Schr\"odinger equation emerges
as its flat-space projection via Fock's 1935 stereographic
map~\cite{Fock1935,loutey_paper7}.

This geometric picture is complete for single atoms.  For molecules, the
current implementation stitches together two separate $S^3$
manifolds---one per atom---using approximate cross terms: Mulliken overlap
for cross-nuclear attraction, Fourier convolution for direct Coulomb
integrals, and STO-overlap bridges for inter-center
hopping~\cite{loutey_fci}.  The result is basis set superposition error
(BSSE), approximate integrals, and no natural coordinate system for the
bonding region.

This paper proposes the resolution: the \emph{bond sphere}.  A diatomic
molecule is described by a single $S^3$ with two weighted poles, one per
nucleus.  The atomic $S^3$ manifolds are the $R \to \infty$ limit.
Cross-atom matrix elements that are currently approximated become exact
$SO(4)$ Wigner $D$-matrix elements.  The natural basis for this geometry
is the Sturmian basis, in which all functions share a single momentum
parameter $p_0$ and therefore live on the same $S^3$ by construction.
The $\beta_k$ eigenvalues that encode nuclear attraction are determined
by the two-center prolate spheroidal eigenvalue problem---whose coordinate
surfaces are precisely the confocal ellipsoids of the bond sphere metric.

The paper is organized as follows.  Sections~\ref{sec:single_atom}
and~\ref{sec:bond_sphere} establish the single-atom and diatomic bond
sphere geometry.  Section~\ref{sec:runge_lenz} gives the Runge--Lenz
alignment argument for why the bond rotation is generated by $A_y$.
Section~\ref{sec:heteronuclear} treats heteronuclear asymmetry---the
egg.  Section~\ref{sec:membrane} identifies the kinetic membrane as
the Fiedler vector zero-crossing.  Section~\ref{sec:sturmian} introduces
the Sturmian basis.  Section~\ref{sec:beta} develops the $\beta_k$
eigenvalue problem via prolate spheroidal separation.
Section~\ref{sec:failure} proves why atom-centered approximations fail.
Section~\ref{sec:uniform} states the uniform diagonal theorem.
Section~\ref{sec:structural_theorem} proves the Sturmian structural
theorem: the Hamiltonian is proportional to the overlap matrix, making
eigenvalues geometry-independent.
Section~\ref{sec:multi_electron} extends the structural theorem to
multi-electron systems via Coulomb and generalized Sturmian CI.
Section~\ref{sec:selection_rules} derives the $SO(4)$ selection rules
for $\sigma$-bond formation.
Section~\ref{sec:conclusion} concludes.


%% ========================================================================
\section{What Is Known vs.\ What Is New}
\label{sec:attribution}
%% ========================================================================

\subsection{Known Results (Prior Work)}

\begin{itemize}

\item \textit{Shibuya and Wulfman (1965)}~\cite{Shibuya1965}: The
two-center Coulomb problem in momentum space maps to a single $S^3$
with two poles.  Sturmian basis functions on this sphere provide a
complete, orthogonal expansion.

\item \textit{Aquilanti and Cavalli
(1980s--1990s)}~\cite{Aquilanti1986,Aquilanti1990}: Hyperspherical
harmonic approach to molecular problems, including systematic use of
$SO(4)$ Wigner $D$-matrices and the relationship between internuclear
distance and angular pole separation on $S^3$.

\item \textit{Avery (2000s)}~\cite{Avery2006}: Extended the Sturmian
basis to general molecular systems with practical computational methods.
De~Fazio, Cavalli, and Aquilanti~\cite{DeFazio2003} developed the
\emph{hyperquantization algorithm}, discretizing scattering problems
on $S^3$ via hyperspherical harmonics---arriving at sparse discrete
Hamiltonians from the continuous side, converging with GeoVac's
graph-first construction (see Sec.~\ref{sec:multi_electron}).
Avery~\cite{Avery2004} formulated the generalized Sturmian secular
equation for multi-electron atoms using configuration-dependent
scaling parameters.

\item \textit{$SO(4)$ representation theory}~\cite{Bander1966,%
Biedenharn1981}: The Wigner $D$-matrix elements of $SO(4)$ in the
$(n, l, m)$ basis, generated by angular momentum $\mathbf{L}$ and the
Runge--Lenz vector $\mathbf{A}$.

\item \textit{Sturmian basis functions}: Established by Shull and
L\"owdin~\cite{shull_lowdin}, Rotenberg~\cite{rotenberg}, and
Aquilanti--Avery~\cite{aquilanti_avery,avery_book}.  Orthogonality
under $1/r$ weight and completeness for the Coulomb problem.

\end{itemize}

\subsection{New Contributions}

\begin{enumerate}

\item \textbf{Bond sphere as a GeoVac lattice object.}  The bond
sphere is the natural molecular generalization of the single-atom
$S^3$ lattice (Papers~0, 1, 7), with within-block edges generated by
$T_\pm, L_\pm$ and cross-block edges weighted by $D$-matrix elements
at $g(\gamma(R))$.

\item \textbf{Runge--Lenz alignment argument}
(Sec.~\ref{sec:runge_lenz}).  The poles are forced onto the bond axis
by dynamics: the Runge--Lenz vector of each atom orients toward the
other nucleus.

\item \textbf{Heteronuclear egg geometry}
(Sec.~\ref{sec:heteronuclear}).  For $Z_A \neq Z_B$, the bond sphere
metric is a prolate spheroid with the heavier nucleus at the fat end.

\item \textbf{Kinetic membrane} (Sec.~\ref{sec:membrane}).  The
Fiedler vector zero-crossing identifies a discrete, combinatorial
analog of Bader's bond critical point.

\item \textbf{Identification that the bond sphere requires the
Sturmian basis} (Sec.~\ref{sec:sturmian}).  The single-$S^3$ geometry
is only consistent when all basis functions share one $p_0$.

\item \textbf{Proof that atom-centered Sturmian bases fail for
heteronuclear molecules} (Sec.~\ref{sec:failure}), with the precise
threshold $\beta > p_0/(2Z_B)$ for binding.

\item \textbf{Uniform diagonal theorem}
(Sec.~\ref{sec:uniform}).  In the molecular Sturmian basis,
$\varepsilon_k = -p_0^2/2$ for every orbital.

\item \textbf{Sturmian structural theorem}
(Sec.~\ref{sec:structural_theorem}).  In the molecular Sturmian basis
with shared $p_0$, the Hamiltonian is proportional to the overlap
matrix; eigenvalues are independent of the bond angle $\gamma$ and
hence of $R$.  Binding requires $\beta_k(R)$ from the full PDE.

\item \textbf{$SO(4)$ selection rules for $\sigma$-bonds}
(Sec.~\ref{sec:selection_rules}).
$D^2_{10,10}(\gamma) = 1$ and $D^2_{00,10}(\gamma) = 0$ for all
$\gamma$: the bond rotation preserves $m = 0$ within $l = 1$ and does
not mix $s$ and $p$ states.

\end{enumerate}


%% ========================================================================
\section{The Single-Atom Case as a Special Limit}
\label{sec:single_atom}
%% ========================================================================

The existing GeoVac single-atom lattice is the $\gamma = 0$ limit of
the bond sphere, so all results of Papers~0--7 are preserved.

\subsection{The $T_\pm, L_\pm$ Operators as $SO(4)$ Generators}

The GeoVac paraboloid lattice is generated by four transition operators
(Paper~7, Eqs.~(5)--(6)):
\begin{align}
L_\pm |n, l, m\rangle &\propto |n, l, m \pm 1\rangle, \\
T_\pm |n, l, m\rangle &\propto |n \pm 1, l, m\rangle.
\end{align}
The angular momentum operators $L_x, L_y, L_z$ satisfy
$\mathfrak{su}(2)$:
\begin{equation}
[L_i, L_j] = i\epsilon_{ijk} L_k.
\label{eq:su2_angular}
\end{equation}
The scaled Runge--Lenz vector $\mathbf{A}$, generated by $T_\pm$
within each $n$-shell, satisfies~\cite{Bander1966}:
\begin{align}
[L_i, A_j] &= i\epsilon_{ijk} A_k, \label{eq:LA_comm} \\
[A_i, A_j] &= i\epsilon_{ijk} L_k. \label{eq:AA_comm}
\end{align}

The $\mathfrak{so}(4) \cong \mathfrak{su}(2) \oplus \mathfrak{su}(2)$
structure is exhibited by
\begin{equation}
\mathbf{J}^{(\pm)} = \tfrac{1}{2}(\mathbf{L} \pm \mathbf{A}),
\label{eq:Jpm_def}
\end{equation}
which satisfy two independent, commuting $\mathfrak{su}(2)$ algebras.
For principal quantum number $n$, both carry spin $j = (n-1)/2$,
giving the $(n-1)/2 \otimes (n-1)/2$ representation with Casimir
\begin{equation}
\mathbf{L}^2 + \mathbf{A}^2 = n^2 - 1,
\label{eq:casimir}
\end{equation}
recovering the $S^3$ eigenvalue spectrum $\lambda_n = -(n^2 - 1)$
of Paper~7.

\subsection{Covariance and the Single-Atom Limit}

Because $T_\pm, L_\pm$ generate $SO(4)$, the graph Laplacian
$L = D - A$ is covariant under $SO(4)$ transformations.  The second
atom can be introduced as an $SO(4)$-rotated copy, with the rotation
angle determined by $R$.  Cross-atom matrix elements are then $SO(4)$
Wigner $D$-matrix elements.

At $\gamma = 0$ (coincident poles): $g(0) = e$, the $D$-matrix reduces
to Kronecker deltas, cross-block edges vanish, and the single-center
lattice of Papers~0--7 is recovered.  The 18 symbolic proofs, the
dimensionless vacuum principle, and the kinetic scale factor $-1/16$
are unchanged.


%% ========================================================================
\section{The Bond Sphere Construction}
\label{sec:bond_sphere}
%% ========================================================================

\subsection{The Angle--Distance Relationship}

The internuclear distance $R$ maps to a polar angle $\gamma$ between
two poles on $S^3$.  The momentum conjugate to $R$ defines
\begin{equation}
p_R = \frac{1}{R},
\label{eq:pR}
\end{equation}
and the angular separation is
\begin{equation}
\cos\gamma = \frac{p_0^2 - p_R^2}{p_0^2 + p_R^2},
\label{eq:cos_gamma}
\end{equation}
which is the fourth component $n_4$ of the Fock projection evaluated
at $|\mathbf{p}| = p_R$.  The sine follows as
\begin{equation}
\sin\gamma = \frac{2p_0\, p_R}{p_0^2 + p_R^2} = \Omega(p_R)\cdot p_R,
\label{eq:sin_gamma}
\end{equation}
where $\Omega$ is the Fock conformal factor.

\paragraph{Limiting behavior.}
At $R \to \infty$: $\gamma \to 0$ (coincident poles, single-atom limit).
At $R \to 0$: $\gamma \to \pi$ (antipodal poles, united-atom limit).

\subsection{Eigenfunctions as $SO(4)$ Wigner $D$-Matrix Elements}

When a second pole is introduced at angle $\gamma$, the cross-center
matrix elements are
\begin{equation}
\braket{n', l', m' | U(g(\gamma)) | n, l, m}
= D^{(n)}_{(l'm'),(lm)}(g(\gamma)),
\label{eq:D_matrix}
\end{equation}
where $g(\gamma)$ is the $SO(4)$ rotation that maps one pole to the
other.  The current Mulliken/Fourier/bridge cross-atom terms are
approximations to these exact $D$-matrix elements.

\subsection{The Bond-Sphere Lattice}

\paragraph{Nodes.}  Each node is labeled $(n, l, m, \alpha)$ with
$\alpha \in \{A, B\}$.  Total spatial orbitals: $2N_{\text{sp}}$.

\paragraph{Within-block edges.}  Generated by $T_\pm, L_\pm$ at each
center, as in the single-atom lattice.

\paragraph{Cross-block edges.}  Weighted by $D$-matrix elements:
\begin{equation}
w_{(n'l'm',B),(nlm,A)} = D^{(n)}_{(l'm'),(lm)}(g(\gamma(R))).
\label{eq:cross_weight}
\end{equation}

\paragraph{Representation-theoretic constraint.}
The $D$-matrix $D^{(n)}_{(l'm'),(lm)}(\gamma)$ is block-diagonal in the
principal quantum number $n$ by the representation theory of $SO(4)$,
not by approximation.  Each $n$-shell corresponds to a distinct irreducible
representation of $SO(4)$ with $j = (n-1)/2$; different irreps are
orthogonal under any group element.  The molecular field (second nucleus at
bond angle $\gamma$) acts as an $SO(4)$ rotation, which cannot mix
inequivalent representations.

Cross-$n$ coupling---for example, Li~$2s$ $\leftrightarrow$ H~$1s$ in
heteronuclear molecules---would require the full $SO(4,2)$ conformal
group, whose generators include dilatations and special conformal
transformations that connect different $n$-shells.  No closed-form
``$SO(4,2)$ $D$-matrix'' for cross-nuclear coupling exists in the
literature, and the standard Shibuya--Wulfman formula is exact only
within a single $n$-manifold.

This block-diagonality is the fundamental reason why LCAO approaches
using $D$-matrix coupling fail to capture $\sigma$ hybridization between
orbitals of different principal quantum number.  The symmetry breaking
$SO(4) \to SO(2)$ caused by the second nucleus
(Sec.~\ref{sec:heteronuclear}) breaks $l$-degeneracy \emph{within} each
$n$-shell but does \emph{not} couple different $n$-shells.

\paragraph{Hamiltonian structure.}
\begin{equation}
H_1 = \begin{pmatrix} H_A & V_{AB} \\ V_{AB}^\dagger & H_B \end{pmatrix},
\label{eq:H1_block}
\end{equation}
where $H_A, H_B$ are single-atom Hamiltonians and $V_{AB}$ has elements
given by Eq.~\eqref{eq:cross_weight} modulated by nuclear charges and
kinetic scale.

\subsection{The Opposite-Sign Rotation}
\label{sec:opposite_sign}

The bond rotation is not a spatial rotation.  Since
$\mathbf{A} = \mathbf{J}^{(+)} - \mathbf{J}^{(-)}$, the rotation
generated by $A_y$ acts as
\begin{equation}
e^{i\gamma A_y}
= e^{i\gamma J^{(+)}_y}\; e^{-i\gamma J^{(-)}_y},
\label{eq:opposite_sign}
\end{equation}
where the $\mathbf{J}^{(+)}$ and $\mathbf{J}^{(-)}$ sectors rotate in
opposite directions: $+\gamma$ and $-\gamma$.  A pure spatial rotation
$e^{i\alpha L_y}$ gives same-sign rotation $d^j(+\alpha) \otimes
d^j(+\alpha)$, preserving radial structure.  The bond rotation gives
$d^j(+\gamma) \otimes d^j(-\gamma)$, mixing $s, p, d, \ldots$
orbitals within each $n$-shell.  This $s$--$p$ hybridization is
precisely what chemical bonding requires.

The $D$-matrix in the $(n, l, m)$ basis is
\begin{multline}
D^{(n)}_{(l'm'),(lm)}(\gamma)
= \sum_{m^+, m^-} \sum_{m'^+, m'^-}
C^{n}_{l'm'; m'^+m'^-}\; C^{n}_{lm; m^+m^-} \\
\times\; d^{j}_{m'^+, m^+}(+\gamma)\;
d^{j}_{m'^-, m^-}(-\gamma),
\label{eq:D_nlm}
\end{multline}
where $j = (n-1)/2$ and $C^{n}_{lm; m^+m^-}$ are the $SO(4)$
Clebsch--Gordan coefficients coupling
$|j, m^+\rangle \otimes |j, m^-\rangle \to |l, m\rangle$ with
$m = m^+ + m^-$.

\subsection{The Antipodal Limit}

At $\gamma = \pi$ ($R = 0$, united atom):
\begin{equation}
D^{(n)}(\pi)^2 = \mathbf{I}_{n^2},
\label{eq:D_pi_squared}
\end{equation}
with eigenvalues $\pm 1$ and parity structure
$|l, m\rangle \to (-1)^{n-1+l+|m|}\, |l, -m\rangle$---the
Runge--Lenz parity, not spatial parity $(-1)^l$.


%% ========================================================================
\section{The Runge--Lenz Alignment Argument}
\label{sec:runge_lenz}
%% ========================================================================

The bond rotation $e^{i\gamma A_y}$ is generated by $A_y$, a component
of the Runge--Lenz vector perpendicular to the bond axis.  This choice
is not a convention but a dynamical consequence.

For a hydrogen atom, the Runge--Lenz vector $\mathbf{A}$ points from
the nucleus toward the perihelion of the classical orbit.  In a
diatomic molecule, the electrostatic field of the partner nucleus
breaks the spherical symmetry of each atom's potential.  The lowest-energy
configuration aligns the Keplerian orbit so that its perihelion points
toward the attracting partner---that is, $\mathbf{A}$ orients along
the internuclear axis.  With $\mathbf{A}$ aligned along $\hat{z}$ (the
bond axis), the rotation that maps one pole to the other is generated
by $A_y$ (or $A_x$), which is perpendicular to the bond.

This is the physical origin of $s$--$p$ hybridization in the bond sphere
picture: the Runge--Lenz vector aligns with the bond, and rotation about
a perpendicular axis mixes orbital angular momentum states.  The
hybridization is not an ansatz imposed on the basis---it is a consequence
of the two-center electrostatic potential acting on the dynamical
symmetry generators of each atom.


%% ========================================================================
\section{Heteronuclear Geometry: The Egg}
\label{sec:heteronuclear}
%% ========================================================================

For a homonuclear diatomic ($Z_A = Z_B$), the bond sphere has uniform
metric and full $SO(4)$ symmetry.  For a heteronuclear diatomic
($Z_A \neq Z_B$), the geometry changes qualitatively.

\subsection{Metric Distortion}

The Fock conformal factor $\Omega = 2p_0/(p^2 + p_0^2)$ maps
momentum space to $S^3$.  When the two nuclei have unequal charges,
the natural momentum scale differs between the two poles: center~$A$
(charge $Z_A$) favors $p_0 \sim Z_A$, while center~$B$ (charge $Z_B$)
favors $p_0 \sim Z_B$.  With a shared $p_0$ (as required by the
single-$S^3$ construction), the conformal factor is distorted
asymmetrically about the two poles.

The result is a bond sphere whose intrinsic metric is not uniform.
The heavier nucleus ($Z_A > Z_B$) occupies a larger solid angle on the
sphere---a ``fat end''---while the lighter nucleus occupies a smaller
solid angle---a ``narrow end.''  The overall shape is a prolate
spheroid, with the heavier nucleus at the fat end and the lighter
nucleus at the narrow end.  We call this the \emph{egg}.

\subsection{Bond Polarity as Egg Asymmetry}

The degree of asymmetry is set by the ratio $Z_A/Z_B$.  For LiH
($Z_A/Z_B = 3$), the egg is noticeably asymmetric.  For H$_2$
($Z_A/Z_B = 1$), the egg is a round sphere.  This asymmetry is
the geometric encoding of bond polarity: a polar bond is an
asymmetric egg, and a nonpolar bond is a symmetric sphere.

\subsection{Topology vs.\ Metric}

The node count per shell ($n^2$ nodes at latitude $n$ from each pole)
is fixed by the Paper~0 packing topology.  The egg deformation changes
the \emph{node weights}---the effective nuclear attraction felt by each
node---but not the \emph{graph topology}.  The adjacency structure
($T_\pm, L_\pm$ edges) is unchanged; only the diagonal elements of the
Hamiltonian (the node weights) and the cross-block edge weights (the
$D$-matrix elements) are affected by the asymmetry.

\subsection{Symmetry Breaking: $SO(4) \to SO(2)$}

The unequal pole weights break $SO(4)$ to the axial $SO(2)$ subgroup
preserving the bond axis.  The magnetic quantum number $m$ remains a
good quantum number (generating $\sigma, \pi, \delta, \ldots$ labels),
while the off-axis generators $L_\pm$ are no longer symmetries:
\begin{align}
[\Delta_{\text{asym}}, L_z] &= 0, \label{eq:Lz_comm} \\
[\Delta_{\text{asym}}, L_\pm] &\neq 0 \quad\text{for } Z_A \neq Z_B.
\label{eq:Lpm_comm}
\end{align}

\subsection{Prolate Spheroidal Coordinates}

The natural coordinate system for the egg is prolate spheroidal
coordinates $(\xi, \eta, \varphi)$ with foci at the two nuclear poles.
The coordinate surfaces are confocal ellipsoids ($\xi = \text{const}$)
and hyperboloids ($\eta = \text{const}$).  These are precisely the
level sets of the distorted bond sphere metric: the ellipsoids trace
curves of constant ``distance from both poles simultaneously,'' while
the hyperboloids trace curves of constant ``distance ratio between
the two poles.''  This coordinate system appears again in
Sec.~\ref{sec:beta} as the natural domain for the $\beta_k$ eigenvalue
problem.


%% ========================================================================
\section{The Kinetic Membrane}
\label{sec:membrane}
%% ========================================================================

Between the two poles of the bond sphere lies a region where the edge
weights from the two centers are comparable in magnitude.  This is the
\emph{kinetic membrane}---not a separate geometric object, but an
emergent feature of the edge weight pattern.

\subsection{Identification via the Fiedler Vector}

In the molecular graph, the kinetic membrane is identified by the
\emph{Fiedler vector}---the eigenvector corresponding to the second
smallest eigenvalue of the graph Laplacian.  The Fiedler vector
partitions the graph into two regions (positive and negative components),
and the membrane is where the Fiedler vector crosses zero.

For a homonuclear diatomic, the Fiedler vector zero-crossing lies at
the midpoint between the two poles, by symmetry.  For a heteronuclear
molecule, the zero-crossing shifts toward the lighter nucleus.  In LiH,
the heavier Li pole ($Z = 3$) dominates a larger volume of the bond
sphere, pushing the membrane toward the H end.

\subsection{Connection to Bader's Bond Critical Point}

The Fiedler vector zero-crossing is the discrete, combinatorial analog
of Bader's bond critical point~\cite{Bader1990} in the Quantum Theory
of Atoms in Molecules (QTAIM).  Bader's bond critical point is defined
as the point where $\nabla\rho = 0$ along the bond path, where $\rho$
is the electron density.  The GeoVac graph gives the same object
without invoking the electron density: the membrane position is
determined purely by the graph Laplacian spectrum.

For heteronuclear molecules, both the Bader critical point and the
Fiedler zero-crossing shift toward the electropositive (lighter) atom.
The membrane is thinner (steeper Fiedler gradient) for strongly polar
bonds and thicker (shallower gradient) for nonpolar bonds.


%% ========================================================================
\section{The Sturmian Basis}
\label{sec:sturmian}
%% ========================================================================

The bond sphere hypothesis requires all basis functions to live on a
single $S^3$ parameterized by one momentum scale $p_0$.  The
hydrogenic eigenstate basis uses $p_0(n) = Z/n$---a different sphere
per shell---which is geometrically inconsistent with the single-$S^3$
construction.  The Sturmian basis resolves this.

\subsection{Definition}

The Sturmian basis functions $S_{nlm}$ are solutions to the scaled
Coulomb problem:
\begin{equation}
\left(-\tfrac{1}{2}\nabla^2 - \frac{\beta_n Z}{r}\right) S_{nlm}
  = E_0\, S_{nlm},
\label{eq:sturmian_eq}
\end{equation}
where $E_0 = -p_0^2/2$ is a fixed reference energy shared by all
functions, and
\begin{equation}
\beta_n = \frac{p_0\, n}{Z}
\label{eq:beta_n}
\end{equation}
is the shell-dependent scaling that enforces this.

The key property: all Sturmian functions share the same $E_0$
regardless of $n$.  Under Fock's projection, all functions live on
the \emph{same} $S^3$ of radius $p_0$.

\subsection{Verification}

Substituting effective charge $\tilde{Z} = \beta_n Z$ into the
hydrogenic eigenvalue:
\begin{equation}
E_0 = -\frac{(\beta_n Z)^2}{2n^2}
     = -\frac{(p_0 n / Z)^2 Z^2}{2n^2}
     = -\frac{p_0^2}{2},
\label{eq:sturmian_verify}
\end{equation}
independent of $n$, as required.

\subsection{Orthogonality under $1/r$ Weight}

Sturmian functions satisfy
\begin{equation}
\left\langle S_{n'l'm'}\left|\frac{1}{r}\right| S_{nlm}\right\rangle
  = \frac{p_0}{n}\,\delta_{nn'}\,\delta_{ll'}\,\delta_{mm'}.
\label{eq:sturmian_orthog}
\end{equation}
The $1/r$-weighted completeness precisely matches the Coulomb kernel
in cross-center integrals, making the two-center problem exact.

\subsection{Geometric Inconsistency of the Hydrogenic Basis}

The hydrogenic eigenstate $\psi_{nlm}$ has $p_0(n) = Z/n$.  For
the cross-center matrix element
$\langle \psi^A_{n'l'm'} | {-Z_B/r_B} | \psi^A_{nlm} \rangle$
with $n' \neq n$, the bra lives on $S^3$ with radius $Z/n'$ and the
ket on $S^3$ with radius $Z/n$.  The integral involves two different
Fock projections---a convolution over mismatched spheres.  This is
the origin of the $\sin\gamma$ form factor in the current SW module:
a lowest-order correction for this geometric mismatch.

In the Sturmian basis, the mismatch is zero by construction.

\subsection{SW Integrals Are Exact in the Sturmian Basis}

This is the central result.  In the Sturmian basis:
\begin{equation}
\boxed{
\left\langle S^A_{n'l'm'}\left| -\frac{Z_B}{r_B}\right|
S^A_{nlm}\right\rangle
= -\frac{Z_B}{p_0}\, D^{(n)}_{(l'm'),(lm)}(\gamma),
}
\label{eq:sw_exact}
\end{equation}
exactly, with no form factor.

The derivation: the cross-center Coulomb potential is related to the
same-center potential by an $SO(4)$ rotation (Eq.~\eqref{eq:D_matrix}).
Using Sturmian orthogonality (Eq.~\eqref{eq:sturmian_orthog}), the
$1/r_A$ matrix element collapses to a Kronecker delta, and the sum
over intermediate states reduces to a single $D$-matrix element.

When the same calculation is attempted with hydrogenic eigenstates of
different $p_0$ values, the orthogonality does not apply, and one
obtains the mismatch integral that, expanded to leading order, gives
the $\sin\gamma$ form factor.

\subsection{The Self-Consistency Condition}

The momentum parameter $p_0$ is determined by
\begin{equation}
p_0 = \sqrt{-2E_{\text{mol}}},
\label{eq:self_consistency}
\end{equation}
where $E_{\text{mol}}$ is the ground-state eigenvalue of $H(p_0)$.
This is a self-consistent fixed point, whose existence follows from
the intermediate value theorem: $f(p_0) = \sqrt{-2E_{\text{mol}}(p_0)}$
maps $\mathbb{R}^+$ to $\mathbb{R}^+$, with $f(p_0) > p_0$ for
small $p_0$ (diffuse basis, weakly negative energy) and $f(p_0) < p_0$
for large $p_0$ (compact basis, bounded energy).

\subsection{$\gamma$ Encodes Both Geometry and Energy}

At self-consistency:
\begin{equation}
\gamma = \arccos\!\left(\frac{p_0^2 - p_R^2}{p_0^2 + p_R^2}\right),
\label{eq:gamma_sc}
\end{equation}
with $p_0 = \sqrt{-2E_{\text{mol}}}$.  Both derivatives
$\partial\gamma/\partial R \neq 0$ and
$\partial\gamma/\partial E_{\text{mol}} \neq 0$, confirming that
$\gamma$ carries information about both bond length and molecular
energy simultaneously.

\subsection{Backward Compatibility with Papers~0--7}

In the limit $p_0 \to Z/n$ per shell: $\beta_n \to 1$, the Sturmian
equation becomes the ordinary hydrogenic equation, and the atomic
eigenvalue $-Z^2/(2n^2)$ is recovered.  The self-consistency loop has
no single-atom analog because $E_n = -Z^2/(2n^2)$ already determines
$p_0(n) = Z/n$.  All single-atom results are the $n$-resolved limit
of the Sturmian construction.


%% ========================================================================
\section{The $\beta_k$ Eigenvalue Problem}
\label{sec:beta}
%% ========================================================================

For a single-center Sturmian, $\beta_n = p_0 n/Z$ is determined by
the nuclear charge alone.  For a two-center molecular Sturmian, $\beta_k$
is the eigenvalue of the prolate spheroidal equation with foci at the
two nuclei.

\subsection{Prolate Spheroidal Separation}

The Sturmian equation with the two-center molecular potential
$V = -Z_A/r_A - Z_B/r_B$ separates in prolate spheroidal coordinates
$(\xi, \eta, \varphi)$ with foci at the nuclear positions.  The
separation yields an angular eigenvalue problem (in $\eta$) and a
radial eigenvalue problem (in $\xi$), coupled by a separation constant
$A$.

The angular equation at fixed magnetic quantum number $m$ and spheroidal
parameter $c = p_0 R/2$ is
\begin{equation}
\frac{d}{d\eta}\!\left[(1-\eta^2)\frac{dS}{d\eta}\right]
+ \left(A_{m,n_{\text{sph}}} - c^2\eta^2
- \frac{m^2}{1-\eta^2}\right)S = 0,
\label{eq:angular_eq}
\end{equation}
whose eigenvalues $A_{m,n_{\text{sph}}}(c, b)$ depend on $c$ and the
charge asymmetry parameter $b = (Z_A - Z_B)R/2$.

The radial equation is
\begin{equation}
\frac{d}{d\xi}\!\left[(\xi^2-1)\frac{dR}{d\xi}\right]
+ \left(-A + c^2\xi^2 + 2a\xi
- \frac{m^2}{\xi^2-1}\right)R = 0,
\label{eq:radial_eq}
\end{equation}
where $a = \beta_k(Z_A + Z_B)R/2$ encodes the molecular Sturmian
eigenvalue $\beta_k$.

\subsection{The Matching Condition}

The molecular $\beta_k$ for each orbital $(m, n_{\text{sph}}, n_{\text{rad}})$
is determined by the matching condition: the angular separation constant
$A_{m,n_{\text{sph}}}(c, b)$ and the radial eigenvalue must be
consistent.  Specifically, $\beta_k$ is found by root-finding (brentq)
on the condition
\begin{equation}
A_{\text{ang}}(c, b) + L_{\text{rad}}(c, a(\beta_k)) = 0,
\label{eq:matching}
\end{equation}
where $L_{\text{rad}}$ is the top radial eigenvalue at the given parameters.

\subsection{Key Result}

\begin{theorem}[Molecular Sturmian $\beta_k$ for LiH]
\label{thm:lih_beta}
For LiH at $R = 3.015$~bohr with $p_0 = \sqrt{-2E_{\text{mol}}}$,
the molecular Sturmian $\beta_k$ values computed via prolate spheroidal
separation are:
\begin{center}
\begin{tabular}{lcc}
Orbital & $\beta_k$ & Character \\
\hline
$1\sigma$ & 1.033 & Li $1s$ core \\
$2\sigma$ & 1.950 & Li-H bonding \\
$1\pi$ & 1.976 & Li-H $\pi$ \\
\end{tabular}
\end{center}
All $n = 3$ orbitals have $\beta_k > 2.3$.  The H$_2^+$ ground-state
energy computed by this method is $E = -1.096$~Ha (exact: $-1.103$~Ha,
0.6\% error).
\end{theorem}

\subsection{Why Prolate Spheroidal Coordinates Are Natural}

The separability of the Sturmian equation in prolate spheroidal
coordinates is not a mathematical accident.  The coordinate foci are
the two nuclear poles, and the coordinate surfaces are the confocal
ellipsoids that are the level sets of the distorted bond sphere metric
(Sec.~\ref{sec:heteronuclear}).  The coordinates respect both poles
simultaneously---which is precisely the requirement for a basis that
lives on the bond sphere rather than on a single-center sphere.


%% ========================================================================
\section{Why Atom-Centered Bases Fail for Heteronuclear Molecules}
\label{sec:failure}
%% ========================================================================

\begin{theorem}[Atom-centered Sturmian failure]
\label{thm:failure}
For a heteronuclear diatomic with $Z_A/Z_B = k > 1$, there is no
atom-centered Sturmian basis with shared $p_0$ that simultaneously
binds both centers.
\end{theorem}

\begin{proof}[Proof sketch.]
The shared $p_0$ is determined by the deeper well ($Z_A$).  At this
$p_0$, the one-electron diagonal for the lightest orbital of center~$B$
is
\begin{equation}
\varepsilon_1(p_0) = \frac{p_0^2}{2} - Z_B\, p_0.
\label{eq:h1s_diag}
\end{equation}
This is positive (unbound) when $p_0 > 2Z_B$.  For the self-consistent
$p_0 \approx \sqrt{-2E_A}$ of the heavier atom with $Z_A \gg Z_B$,
this condition is generically satisfied.

The cross-nuclear correction from $Z_A$ on the $B$-center $1s$ function,
evaluated in the atom-centered Sturmian basis, is
$\Delta\varepsilon = -(Z_A/p_0) D^{(1)}(\gamma)$.  For $n = 1$,
$D^{(1)} = 1$, so $\Delta\varepsilon = -Z_A/p_0$.  The total diagonal
becomes
\begin{equation}
\varepsilon_1^{\text{tot}} = \frac{p_0^2}{2} - Z_B p_0 - \frac{Z_A}{p_0}.
\label{eq:h1s_total}
\end{equation}
For binding: $\varepsilon_1^{\text{tot}} < 0$, i.e.,
$\beta_{\text{eff}} > p_0/(2Z_B)$.  But the atom-centered $\beta_1 = p_0/Z_B$
gives $\varepsilon_1 = p_0^2/2 - Z_B p_0$ plus the cross-nuclear
$-Z_A/p_0$.  At $p_0 = 3.168$ (LiH): $\varepsilon_1^{\text{tot}} =
5.02 - 3.17 - 0.95 = +0.85$~Ha---still positive.

The cross-nuclear correction is insufficient because the atom-centered
function has the wrong spatial extent: it was defined against $Z_B$
alone and is too compact near the $B$ center to capture the full
$A$-$B$ cross-nuclear attraction.  The correct integral requires the
\emph{molecular} Sturmian function---the two-center prolate spheroidal
eigenfunction---which has the correct spatial extent near both centers.
\end{proof}

\begin{corollary}
For LiH ($Z_A/Z_B = 3$): the binding threshold is
$\beta > p_0/(2Z_B) = 1.581$.  The $2\sigma$ molecular Sturmian
achieves $\beta = 1.950$ (Theorem~\ref{thm:lih_beta}), but no
atom-centered approximation reaches this threshold.
\end{corollary}

This argument is independent of the number of electrons and applies to
any heteronuclear diatomic with $Z_A/Z_B$ sufficiently different from~1.
The diagnostic arc v0.9.20--v0.9.30 (documented in the companion
paper~\cite{loutey_fci}) exhaustively verifies this theorem
computationally: every atom-centered variant fails for LiH by the
mechanism described above.


%% ========================================================================
\section{The Uniform Diagonal Theorem}
\label{sec:uniform}
%% ========================================================================

\begin{theorem}[Uniform diagonal]
\label{thm:uniform}
In the molecular Sturmian basis with shared $p_0$, the diagonal of the
one-electron Hamiltonian is
\begin{equation}
\varepsilon_k = -\frac{p_0^2}{2}
\label{eq:uniform_diag}
\end{equation}
for every orbital $k$, regardless of quantum numbers.
\end{theorem}

\begin{proof}
Each molecular Sturmian $S_k$ satisfies
$(-\frac{1}{2}\nabla^2 - \beta_k V_{\text{mol}}) S_k = E_0 S_k$
with $E_0 = -p_0^2/2$.  The diagonal matrix element of the full
one-electron Hamiltonian $\hat{h}_1 = -\frac{1}{2}\nabla^2 + V_{\text{mol}}$
is
\begin{equation}
\langle S_k | \hat{h}_1 | S_k \rangle
= E_0 + (1 - \beta_k^{-1})\,
\langle S_k | V_{\text{mol}} | S_k \rangle.
\end{equation}
In the exact molecular Sturmian basis with proper normalization, the
$V_{\text{mol}}$ expectation value is $-p_0^2/\beta_k$, giving
\begin{equation}
\varepsilon_k = -\frac{p_0^2}{2} - \frac{p_0^2}{\beta_k}
+ \frac{p_0^2}{\beta_k} = -\frac{p_0^2}{2}.
\end{equation}
\end{proof}

\paragraph{Physical interpretation.}
The uniform diagonal means that no orbital has a different diagonal
energy.  The entire information content of the nuclear attraction is
in the \emph{off-diagonal} $\beta_k$ structure and the $D$-matrix
coupling.  The asymmetry between the two centers enters only through
the $\beta_k$ eigenvalues (which differ for $\sigma$ vs.\ $\pi$ MOs
and depend on $Z_A, Z_B, R$) and through the $D$-matrix rotation
angle $\gamma$.

\paragraph{Consistency with the $\beta_k$ discussion.}
The $\beta_k$ eigenvalues (Sec.~\ref{sec:beta}) parametrize the
nuclear attraction: orbital $k$ has effective charge $\beta_k Z$
in the Sturmian equation.  This information enters the FCI Hamiltonian
through the off-diagonal matrix elements (the $D$-matrix coupling
weighted by $\beta_k$-dependent factors), not through the diagonal.
The diagonal is uniformly $-p_0^2/2$, and all orbital differentiation
is off-diagonal.  This is the correct architecture for a bond sphere
FCI: the basis is adapted to the molecular potential, and the nuclear
charges appear in the coupling structure, not in the orbital energies.


%% ========================================================================
\section{Geometric Limits of the Sturmian Bond Sphere}
\label{sec:structural_theorem}
%% ========================================================================

The computational diagnostic arc v0.9.20--v0.9.34 (29~versions) exhaustively
tested the Sturmian bond sphere for LiH.  Every variant failed.  This section
distills the failure into a single structural theorem that explains why the
Sturmian basis, despite its geometric elegance, cannot produce
geometry-dependent binding in the shared-$p_0$ framework without solving a
continuous PDE at every~$R$.

\begin{theorem}[Sturmian structural theorem]
\label{thm:structural}
In the molecular Sturmian basis with shared momentum scale $p_0$,
the one-electron Hamiltonian matrix satisfies
\begin{equation}
H_{ij} = \frac{p_0^2}{2}\!\left(\frac{1}{\beta_j} - 1\right) S_{ij},
\label{eq:H_prop_S}
\end{equation}
where $S_{ij}$ is the overlap matrix and $\beta_j$ is the Sturmian
eigenvalue of orbital~$j$.  The Hamiltonian is proportional to the overlap
matrix.  Consequently, the generalized eigenvalues
\begin{equation}
E_k = \frac{p_0^2}{2}\!\left(\frac{1}{\beta_k} - 1\right)
\label{eq:E_geom_indep}
\end{equation}
are independent of the bond angle $\gamma$ and hence independent of the
internuclear distance~$R$.
\end{theorem}

\begin{proof}
Each Sturmian function $S_j$ satisfies
$(-\tfrac{1}{2}\nabla^2 - \beta_j V_{\text{mol}}) S_j = E_0 S_j$
with $E_0 = -p_0^2/2$.  The matrix element of the physical Hamiltonian
$\hat{h}_1 = -\tfrac{1}{2}\nabla^2 + V_{\text{mol}}$ between Sturmians
$S_i$ and $S_j$ is
\begin{align}
H_{ij} &= \langle S_i | \hat{h}_1 | S_j \rangle \notag \\
       &= E_0\, S_{ij} + (1 - \beta_j^{-1})
          \langle S_i | V_{\text{mol}} | S_j \rangle.
\end{align}
But $\langle S_i | V_{\text{mol}} | S_j \rangle
= -\beta_j^{-1}\langle S_i | (\hat{h}_1 - E_0) | S_j \rangle
= -\beta_j^{-1}(H_{ij} - E_0 S_{ij})$, which gives the recursion.
When $\beta_j$ is treated as a fixed parameter (independent of $R$),
the overlap matrix $S_{ij}$ depends on $\gamma(R)$ through the
$D$-matrix coupling, but $H_{ij}/S_{ij}$ does not.  The generalized
eigenvalue problem $H\mathbf{c} = E\, S\mathbf{c}$ then yields
eigenvalues $E_k$ that depend only on $\{p_0, \beta_k\}$, not on
$\gamma$.  The nuclear repulsion $V_{NN} = Z_A Z_B / R$ adds a purely
$1/R$ potential, producing a monotonically repulsive total energy with
no minimum.  This was confirmed numerically in v0.9.34 diagnostic
(bond sphere Hamiltonian analysis): $E_{\text{total}}(R)$ is a
pure~$1/R$ curve with no bound state.
\end{proof}

\begin{corollary}[Binding requires $\beta_k(R)$]
\label{cor:binding}
Binding in the Sturmian framework requires the $R$-dependent prolate
spheroidal eigenvalue spectrum $\beta_k(R)$.  This in turn requires
solving the two-center Schr\"odinger PDE at every~$R$, which
reintroduces continuous spatial geometry and violates the discrete graph
premise of GeoVac.
\end{corollary}

\begin{corollary}[Heteronuclear dual-$p_0$ failure]
\label{cor:dual_p0}
For heteronuclear diatomics, no shared $p_0$ exists that simultaneously
represents both atomic length scales.  The dual-$p_0$ extension
(v0.9.34) fails because molecular Sturmian functions are inherently
delocalized over both centers: at $p_{0,H} = 1.0$, the H-scale
$1\sigma'$ orbital ``sees'' the Li nucleus ($Z_A = 3$) and acquires
$H_1 = -2.6$~Ha.  The self-consistent loop drives both $p_0$ values
toward a single geometric mean ($p_0 \approx 2.0$), eliminating the
scale separation that was the motivation for the extension.
\end{corollary}

\paragraph{Sturmian basis for heteronuclear systems.}
The Sturmian basis resolves the geometric mismatch between $n$-shells by
placing all orbitals on a single $S^3$ with shared momentum scale $p_0$.
In this basis, different Sturmian $n$-shells genuinely do couple through
the cross-nuclear operator $-Z_B/r_B$ (Eq.~\eqref{eq:sw_exact}), making
the $D$-matrix formula exact rather than approximate---in contrast to the
hydrogenic basis, where the block-diagonality constraint forbids cross-$n$
coupling.

However, the Sturmian approach fails for heteronuclear molecules (e.g.,
LiH) because no single $p_0$ simultaneously satisfies the
self-consistency condition $p_0 = \sqrt{-2E}$ for both atoms
(Corollary~\ref{cor:dual_p0}).  Twenty-nine diagnostic versions
(v0.9.20--34) explored Sturmian-based molecular Hamiltonians; all
exhibited pathological behavior for LiH due to this $p_0$
incompatibility.  The heteronuclear problem thus remains open.


%% ========================================================================
\section{Multi-Electron Sturmian CI: Computational Verification}
\label{sec:multi_electron}
%% ========================================================================

The structural theorem (Theorem~\ref{thm:structural}) proves
$H \propto S$ for one-electron molecular Sturmians with shared
momentum scale~$p_0$.  The natural question is whether multi-electron
Sturmians---where the electron--electron repulsion $V_{ee}$ introduces
coupling beyond one-body structure---can circumvent this limitation and
produce basis sets that outperform standard hydrogenic FCI for atoms
and molecules.  Tracks~BU-1 and BU-1b (v2.0.33) tested this for
helium (two electrons) using both Coulomb and generalized Sturmian CI.

\paragraph{Convergent discovery: Aquilanti and Avery.}
The GeoVac framework and the Aquilanti group
(Perugia)~\cite{Aquilanti1986,Aquilanti1990,Aquilanti2001} independently
arrived at the same mathematical structure---Fock projection onto~$S^3$,
hyperspherical harmonic expansions, angular momentum algebra via
Wigner~$3j$ and Gaunt coefficients, and sparse discrete
Hamiltonians---from opposite directions.  GeoVac proceeds graph $\to$
continuous (Paper~7~\cite{loutey_paper7}): the discrete graph Laplacian
is the primary object, and the Schr\"odinger equation emerges via
stereographic projection.  Aquilanti and collaborators proceed continuous
$\to$ discrete: the hyperquantization
algorithm~\cite{DeFazio2003,AquilantiCapecchi2000} discretizes continuous
scattering problems on~$S^3$ via harmonic analysis and discrete
polynomials.  Avery~\cite{Avery2004} formulated the generalized Sturmian
secular equation for multi-electron systems, using
configuration-dependent scaling parameters that are closely related to
the $\beta_k$ eigenvalues of Sec.~\ref{sec:beta}.  Neither group
previously combined these bases with qubit encoding for quantum
simulation---the application explored in
Paper~14~\cite{loutey_paper14}.

\paragraph{Methods.}
Two Sturmian variants were tested against standard hydrogenic FCI for
the He ground state (exact: $-2.9037$~Ha).

\emph{Coulomb Sturmian CI:} Each orbital $(n,l,m)$ uses
$Z_{\mathrm{eff}} = n k$, where $k$ is a common scaling parameter
optimized variationally.  The basis is L\"owdin-orthogonalized before
standard FCI via Slater--Condon rules with full $R^k$ Slater integrals
and Gaunt angular coefficients.

\emph{Generalized Sturmian CI:} Each Slater determinant~$\nu$ receives
its own scaling $\beta_\nu$ from the isoenergetic condition
$E_{\mathrm{trial}} = -\beta_\nu^2 Z^2 / 2 \sum_i 1/n_i^2$.
All orbitals within one configuration share
$Z_{\mathrm{eff}} = \beta_\nu Z$.  A self-consistency loop iterates
$E_{\mathrm{trial}} \to \beta_\nu \to (H, S) \to$ generalized
eigenvalue $\to E_{\mathrm{new}}$ until convergence.

\paragraph{Results.}

\begin{table}[h]
\caption{He ground state energy comparison (exact: $-2.9037$~Ha).
  $\Delta E$ is energy difference from standard FCI at same
  $n_{\max}$.}
\label{tab:sturmian_he}
\begin{ruledtabular}
\begin{tabular}{llccr}
$n_{\max}$ & Method & $E$ (Ha) & Error (\%) & $\Delta E$ (mHa) \\
\hline
2 & Standard FCI       & $-2.8078$ & 3.30 & --- \\
2 & Coulomb Sturmian   & $-2.8116$ & 3.17 & $-3.8$ \\
2 & Generalized        & $-2.8249$ & 2.72 & $-17.1$ \\
\hline
3 & Standard FCI       & $-2.8172$ & 2.98 & --- \\
3 & Coulomb Sturmian   & $-2.8440$ & 2.06 & $-26.8$ \\
3 & Generalized        & $-2.8319$ & 2.47 & $-14.7$ \\
\end{tabular}
\end{ruledtabular}
\end{table}

The Coulomb Sturmian improves on standard FCI at both basis sizes,
with the advantage growing from 0.13~percentage points (pp) at
$n_{\max} = 2$ to 0.92~pp at $n_{\max} = 3$.  The generalized
Sturmian is best at $n_{\max} = 2$ ($-17.1$~mHa, 2.72\% error) but
exhibits crossover at $n_{\max} = 3$: it is 12.1~mHa \emph{worse}
than the Coulomb variant.

\paragraph{$\beta$ differentiation.}
The generalized Sturmian at $n_{\max} = 3$ produces six distinct
$\beta$ values ranging from $0.84$ (the $1s^2$ configuration) to
$2.52$ (the $3d^2$ configuration)---a factor~$3\times$ spread.
Configuration-level differentiation \emph{is} present: compact core
configurations receive smaller~$\beta$ (more diffuse orbitals), while
high-$n$ configurations receive larger~$\beta$ (more contracted).
However, the isoenergetic constraint forces all orbitals within a
single configuration to share one $Z_{\mathrm{eff}} = \beta_\nu Z$,
removing the $n$-dependent radial differentiation that the Coulomb
Sturmian provides ($Z_{\mathrm{eff}} = nk$).  At small basis
($n_{\max} = 2$), between-configuration flexibility dominates and the
generalized method wins.  At larger basis ($n_{\max} = 3$), the loss
of within-configuration orbital flexibility outweighs this gain.

\paragraph{Structural finding.}
The generalized Sturmian's isoenergetic constraint removes
$n$-dependent radial differentiation---a multi-electron analog of the
single-electron structural theorem
(Theorem~\ref{thm:structural}).  In the one-electron case,
$H \propto S$ makes eigenvalues geometry-independent.  In the
multi-electron case, the constraint
$Z_{\mathrm{eff}} = \beta_\nu Z$ (uniform within each
configuration) limits the radial flexibility that $V_{ee}$ can
exploit, and no variant tested overcomes the fundamental tradeoff
between mathematical elegance and radial flexibility.

\paragraph{Gaunt sparsity preservation.}
The Gaunt selection rule sparsity is exactly preserved in both
Sturmian variants: the electron repulsion integral (ERI) nonzero
counts are identical to the standard hydrogenic basis at each
$n_{\max}$ (65~at $n_{\max} = 2$; 1492~at $n_{\max} = 3$).
This is expected: Gaunt coefficients depend only on
$(l, m)$ quantum numbers, not on radial functions or scaling
parameters.  The result confirms that the angular selection rule
structure identified in Paper~7~\cite{loutey_paper7} is
\emph{basis-independent}---a property of the $S^3$ topology, not of
the specific radial representation.

\paragraph{Qubit encoding comparison.}
Track~BU-2 compared the Sturmian and standard qubit Hamiltonians
(via L\"owdin orthogonalization and Jordan--Wigner transformation;
see Paper~14~\cite{loutey_paper14} for methodology).
At $Q = 10$ qubits ($n_{\max} = 2$), the Sturmian Hamiltonian has
$\sim$7\% fewer Pauli terms (112 vs.\ 120) but $2.8\times$ higher
Pauli 1-norm (31.4 vs.\ 11.3~Ha).  At $Q = 28$ ($n_{\max} = 3$),
the Pauli term count is nearly identical (2627 vs.\ 2659), while
the 1-norm inflation grows to $4.5\times$ (349 vs.\ 78~Ha).
The 1-norm inflation is structural: the L\"owdin $S^{-1/2}$
transformation spreads overlap matrix structure into coefficient
magnitudes, inflating operator weights throughout the qubit
Hamiltonian.  Since the Pauli 1-norm controls Trotter error bounds
and measurement shot budgets, the Sturmian basis is
\emph{disadvantaged} for quantum simulation despite its modest
energetic advantages in classical FCI.

\paragraph{Negative conclusion.}
The multi-electron Sturmian direction does not bypass the accuracy
ceilings established by the natural geometry hierarchy
(Papers~11--17).  The investigation yields four structural findings:
(a)~the single-electron structural theorem extends qualitatively to
multi-electron systems---structural constraints on scaling parameters
limit radial flexibility regardless of $V_{ee}$ coupling;
(b)~the convergent discovery with Aquilanti and Avery confirms that
$S^3$ harmonic analysis and Gaunt-based angular algebra are
canonical structures, arrived at independently from both directions;
(c)~Gaunt selection rule sparsity is basis-independent, depending only
on the angular quantum numbers, not on radial functions;
(d)~L\"owdin orthogonalization inflates the Pauli 1-norm by
$3$--$5\times$, making nonorthogonal bases structurally disadvantaged
for qubit Hamiltonians.


%% ========================================================================
\section{$SO(4)$ Selection Rules for $\sigma$-Bond Formation}
\label{sec:selection_rules}
%% ========================================================================

The $SO(4)$ Wigner $D$-matrix encodes the full angular coupling between
atomic orbitals under the bond rotation $e^{i\gamma A_y}$.  Two diagonal
elements exhibit remarkable constancy.

\begin{theorem}[$\sigma$-bond selection rules]
\label{thm:selection}
For the $n = 2$ representation of $SO(4)$:
\begin{align}
D^{(2)}_{(1,0),(1,0)}(\gamma) &= 1 \quad \text{for all } \gamma,
\label{eq:D2_1010} \\
D^{(2)}_{(0,0),(1,0)}(\gamma) &= 0 \quad \text{for all } \gamma.
\label{eq:D2_0010}
\end{align}
That is, the bond rotation generated by $A_y$:
\begin{enumerate}
\item preserves the $2p_0$ orbital ($l = 1, m = 0$) without any
attenuation at any bond angle, and
\item does not mix the $2s$ orbital ($l = 0, m = 0$) with the $2p_0$
orbital ($l = 1, m = 0$).
\end{enumerate}
\end{theorem}

\begin{proof}
The $D$-matrix in the $(l, m)$ basis is computed from the
$\mathbf{J}^{(\pm)}$ Clebsch--Gordan decomposition
(Eq.~\eqref{eq:D_nlm}) with $j = (n-1)/2 = 1/2$ for $n = 2$.
The opposite-sign rotation $d^{1/2}(+\gamma) \otimes d^{1/2}(-\gamma)$
yields, for the $(l' = 1, m' = 0) \to (l = 1, m = 0)$ element:
\begin{equation}
D^{(2)}_{(1,0),(1,0)} = \cos^2(\gamma/2) + \sin^2(\gamma/2) = 1.
\end{equation}
For the $(l' = 0, m' = 0) \to (l = 1, m = 0)$ element:
\begin{equation}
D^{(2)}_{(0,0),(1,0)} = \cos(\gamma/2)\sin(\gamma/2)
                       - \sin(\gamma/2)\cos(\gamma/2) = 0.
\end{equation}
Both identities hold for all $\gamma \in [0, \pi]$.
\end{proof}

\paragraph{Physical interpretation.}
The first identity states that a $2p_0$ ($\sigma$) orbital is perfectly
transmitted through the bond rotation at any internuclear distance---it is
a \emph{transparent mode} of the $SO(4)$ bond channel.  The second states
that the $2s$ orbital cannot scatter into $2p_0$ via the bond rotation
alone.  Together, these enforce the quantum mechanical selection rule that
$\sigma$-bonds are formed by $p_z$ orbitals aligned along the bond axis,
not by $s$--$p$ mixing through the Runge--Lenz rotation.  The $s$--$p$
hybridization that creates $\sigma$-bonds in the conventional picture
arises from the \emph{energy optimization} (variational principle), not
from the geometric coupling.

\paragraph{Numerical verification.}
These identities were verified numerically across $10{,}000$ values of
$\gamma \in [0.01, \pi - 0.01]$ using the \texttt{wigner\_D\_so4}
implementation (v0.9.15).  The computed values satisfy
$|D^{(2)}_{(1,0),(1,0)}(\gamma) - 1| < 10^{-14}$ and
$|D^{(2)}_{(0,0),(1,0)}(\gamma)| < 10^{-14}$ at every grid point
(\texttt{debug/test\_harmonic\_phase\_lock.py}).


%% ========================================================================
\section{Conclusion}
\label{sec:conclusion}
%% ========================================================================

This paper establishes two classes of results for the bond sphere---the
description of a diatomic molecule as a single $S^3$ with two weighted
poles.

\paragraph{Positive structural results.}
The bond sphere geometry is consistent and theoretically motivated.
The internuclear distance encodes as a polar angle $\gamma$ via
$\cos\gamma = (p_0^2 - p_R^2)/(p_0^2 + p_R^2)$, and all cross-atom
matrix elements are $SO(4)$ Wigner $D$-matrix elements evaluated at the
bond rotation $e^{i\gamma A_y}$ (Eq.~\eqref{eq:D_matrix}).  The
Runge--Lenz alignment argument (Sec.~\ref{sec:runge_lenz}) shows that
$s$--$p$ hybridization is a dynamical consequence of the two-center
potential.  For heteronuclear molecules, the bond sphere becomes an
egg---a prolate spheroid whose asymmetry encodes bond polarity
(Sec.~\ref{sec:heteronuclear}).  The kinetic membrane
(Sec.~\ref{sec:membrane}) provides a discrete analog of Bader's bond
critical point.  The $SO(4)$ selection rules
(Theorem~\ref{thm:selection}) show that $D^2_{10,10}(\gamma) = 1$ and
$D^2_{00,10}(\gamma) = 0$ for all $\gamma$: the $2p_0$ orbital is a
transparent mode of the bond channel, and $s$--$p$ mixing requires the
variational principle, not the geometric coupling.  The single-atom
lattice of Papers~0--7 is the $\gamma = 0$ limit; the 18~symbolic
proofs and the kinetic scale factor $-1/16$ are unchanged.

\paragraph{Negative theorem: Sturmian basis limits.}
The Sturmian structural theorem (Theorem~\ref{thm:structural}) proves
that in the shared-$p_0$ Sturmian basis, the one-electron Hamiltonian is
proportional to the overlap matrix, making eigenvalues independent of the
bond angle $\gamma$ and hence of~$R$.  Binding requires the full
$R$-dependent prolate spheroidal eigenvalue spectrum $\beta_k(R)$, which
reintroduces continuous spatial geometry (Corollary~\ref{cor:binding}).
For heteronuclear diatomics, no shared $p_0$ simultaneously represents
both atomic length scales (Corollary~\ref{cor:dual_p0}).  This theorem
was verified computationally across 29~diagnostic versions
(v0.9.20--v0.9.34), all of which failed for LiH by the predicted
mechanism.

\paragraph{Computational path forward.}
The bond sphere is geometrically consistent but its computational
implementation requires either the full LCAO architecture---which
preserves atom-specific length scales and produces $D_e^{\text{CP}}
= 0.093$~Ha at nmax$= 3$ (1.0\% error vs.\ experiment, documented in
the companion paper~\cite{loutey_fci})---or a non-Sturmian basis that
avoids the proportionality constraint of Theorem~\ref{thm:structural}.
The LCAO path is validated; the bond sphere geometry provides the
theoretical framework within which the LCAO results are understood.

\paragraph{Resolution: the prolate spheroidal lattice.}
The failure of the single-$p_0$ Bond Sphere approach points toward the
correct solution: rather than forcing both centers onto a shared $S^3$
with one momentum parameter, each center retains its identity while
being embedded in the natural two-center geometry---the prolate
spheroid.  In prolate spheroidal coordinates $(\xi, \eta, \varphi)$,
the two-center Coulomb problem separates exactly, and discretizing
the Laplace-Beltrami operator on this manifold yields a molecular
lattice with $R$-dependence built into the metric.  For H$_2^+$, this
approach achieves $R_{\mathrm{eq}} = 2.001$~bohr (0.21\% error) and
$E_{\min} = -0.598$~Ha (0.70\% error) with zero fitted
parameters~\cite{loutey_paper11}.


%% ========================================================================
\begin{acknowledgments}
Computational resources were provided by personal hardware.
\end{acknowledgments}
%% ========================================================================

\begin{thebibliography}{25}

\bibitem{loutey_paper0}
J.~Loutey,
``The Geometric Atom: Quantum State Space as a Packing Problem,''
GeoVac Paper~0 (2026).

\bibitem{loutey_paper1}
J.~Loutey,
``The Geometric Atom: Quantum Mechanics as a Packing Problem,''
GeoVac Paper~1 (2026).

\bibitem{loutey_paper7}
J.~Loutey,
``The Dimensionless Vacuum: Recovering the Schr\"odinger Equation
from Scale-Invariant Graph Topology,''
GeoVac Paper~7 (2026).

\bibitem{loutey_fci}
J.~Loutey,
``Full configuration interaction on a sparse graph Hamiltonian:
Multi-electron atoms via relational lattice indexing,''
GeoVac FCI Paper (2026).

\bibitem{Fock1935}
V.~Fock,
``Zur Theorie des Wasserstoffatoms,''
Z. Phys. \textbf{98}, 145--154 (1935).

\bibitem{Shibuya1965}
T.~Shibuya and C.~E. Wulfman,
``Molecular orbitals in momentum space,''
Proc. R. Soc. Lond. A \textbf{286}, 376--389 (1965).

\bibitem{Aquilanti1986}
V.~Aquilanti, S.~Cavalli, and G.~Grossi,
``Hyperspherical coordinates for molecular dynamics by the method
of trees and the mapping of potential energy surfaces for triatomic
systems,''
J. Chem. Phys. \textbf{85}, 1362--1375 (1986).

\bibitem{Aquilanti1990}
V.~Aquilanti, S.~Cavalli, C.~Coletti, and G.~Grossi,
``Alternative Sturmian bases and momentum space applications,''
Chem. Phys. \textbf{209}, 405--416 (1996).

\bibitem{Aquilanti2001}
V.~Aquilanti, A.~Caligiana, S.~Cavalli, and C.~Coletti,
``Hydrogenic orbitals in momentum space and hyperspherical harmonics:
elliptic Sturmian basis sets,''
Int. J. Quantum Chem. \textbf{92}, 212--228 (2003).

\bibitem{Avery2006}
J.~Avery,
\textit{Hyperspherical Harmonics and Generalized Sturmians}
(Kluwer, Dordrecht, 2000).

\bibitem{Bander1966}
M.~Bander and C.~Itzykson,
``Group theory and the hydrogen atom (I),''
Rev. Mod. Phys. \textbf{38}, 330--345 (1966).

\bibitem{Biedenharn1981}
L.~C. Biedenharn and J.~D. Louck,
\textit{Angular Momentum in Quantum Physics},
Encyclopedia of Mathematics and its Applications, Vol.~8
(Addison-Wesley, Reading, MA, 1981).

\bibitem{shull_lowdin}
H.~Shull and P.-O.~L\"owdin,
``Superposition of configurations and natural spin orbitals.
Applications to the He problem,''
J. Chem. Phys. \textbf{25}, 1035 (1956).

\bibitem{rotenberg}
M.~Rotenberg,
``Theory and application of Sturmian functions,''
Adv. At. Mol. Phys. \textbf{6}, 233--268 (1970).

\bibitem{aquilanti_avery}
V.~Aquilanti and J.~Avery,
``Generalized potential harmonics and contracted Sturmians,''
Chem. Phys. Lett. \textbf{267}, 1--8 (1997).

\bibitem{avery_book}
J.~Avery,
\textit{Hyperspherical Harmonics and Generalized Sturmians}
(Kluwer, Dordrecht, 2000).

\bibitem{Bader1990}
R.~F.~W.~Bader,
\textit{Atoms in Molecules: A Quantum Theory}
(Oxford University Press, Oxford, 1990).

\bibitem{Boys1970}
S.~F. Boys and F.~Bernardi,
``The calculation of small molecular interactions by the differences
of separate total energies. Some procedures with reduced errors,''
Mol. Phys. \textbf{19}, 553--566 (1970).

\bibitem{loutey_paper11}
J.~Loutey,
``The Molecular Fock Projection: Diatomic Molecules as Prolate
Spheroidal Lattices,''
GeoVac Paper~11 (2026).

\bibitem{loutey_paper14}
J.~Loutey,
``Structurally Sparse Qubit Hamiltonians from Natural Geometry,''
GeoVac Paper~14 (2026).

\bibitem{DeFazio2003}
P.~De~Fazio, S.~Cavalli, and V.~Aquilanti,
``Orthogonal polynomials of a discrete variable as expansion basis
sets in quantum mechanics: Hyperquantization algorithm,''
Int. J. Quantum Chem. \textbf{93}, 91--111 (2003).

\bibitem{AquilantiCapecchi2000}
V.~Aquilanti and G.~Capecchi,
``Harmonic analysis and discrete polynomials. From semiclassical
information to time-frequency analysis,''
Theor. Chem. Acc. \textbf{104}, 183--188 (2000).

\bibitem{Avery2004}
J.~Avery,
``The generalized Sturmian method and inelastic scattering of fast
electrons,''
J. Phys. Chem. A \textbf{108}, 8848--8851 (2004).

\end{thebibliography}

\end{document}
